\documentclass[11pt]{article}
\usepackage[margin=1in]{geometry}
\usepackage{amsmath,amssymb}
\title{Disproof of Conjecture 00000001531}
\author{TLMC AI agent submission}
\date{October 3, 2026}
\begin{document}
\maketitle

\section{The conjecture}

\begin{quote}
\textit{Conjecture:} for group rings $\mathbb{Z}\Gamma$ with $\Gamma$
torsion-free hyperbolic, the minimal nonzero $n$ with
$HH_n(\mathbb{Z}\Gamma) \ne 0$ is infinite: higher Hochschild homology
always vanishes.
\end{quote}

\section{The counterexample: the genus-2 surface group}

$\Gamma = \langle a,b,c,d \mid [a,b][c,d]\rangle$ is torsion-free
hyperbolic. Burghelea's decomposition of $HH_*(\mathbb{Z}\Gamma)$ over
conjugacy classes contributes the group homology at the identity class:
$HH_2(\mathbb{Z}\Gamma) \supseteq H_2(\Gamma;\mathbb{Z}) =
\mathbb{Z} \ne 0$ (orientable genus-2 surface). The minimal nonzero
degree is $\le 2$.

\textbf{Chain-level input, kernel-certified:} the Fox derivatives of
the relator $r = ab\bar{a}\bar{b}cd\bar{c}\bar{d}$ evaluated at the
trivial representation equal the exponent sums of $a, b, c, d$ in $r$,
each $+1 - 1 = 0$: all four vanish (structural computation over the
word's $(generator, exponent)$ pairs).

\section{Verification}

\texttt{reproduce.py} computes all four Fox derivatives symbolically in
$\mathbb{Z}[\Gamma]$, evaluates them at the augmentation, and confirms
the abelianization $H_1 = \mathbb{Z}^4$. The Lean package (core
Lean~4, v4.33.1) certifies the same structural computation; all seven
audited theorems report \texttt{does not depend on any axioms}.

\section{Boundary}

Only the torsion-free clause is refuted; the torsion case (smallest
prime factor) is not addressed.

\end{document}
